\subsection{Interpretability of the Learned Parameters}
\label{app:interpretability}

In Red-LGCP, whether a cell is active is decided by the classifier, a neural network whose weights have no direct interpretation, whereas the counts of active cells follow the LGCP intensity $\log\lambda_i=\alpha+\mathbf{x}_i^\top\boldsymbol{\beta}+f(\mathbf{r}_i)$, whose parameters do. Table~\ref{tab:lgcp_params} reports them as mean $\pm$ standard deviation over the 10 seeds $\times$ 4 weekly windows of each month. Since the covariates are standardized, we convert each coefficient $\beta_j$ into the multiplicative effect $\exp(\beta_j/\mathrm{sd}_j)$ on $\lambda$ of a one-unit increase of the raw covariate, or of switching a binary indicator on, where $\mathrm{sd}_j$ is the training standard deviation of covariate $j$ in each window (for chlorophyll-a, we report $\exp(\beta_j)$, the effect of one standard deviation).

\begin{table}[t]
\centering
\caption{Parameters of the LGCP intensity of Red-LGCP (mean $\pm$ population standard deviation over 10 seeds $\times$ 4 weekly windows per month). The effect columns give the multiplicative effect on the intensity $\lambda$ of active cells, computed per run with the training statistics of its window. Since the kernel is a product of a spatial and a temporal kernel, the spatial and the temporal variances are identified only up to a common factor. The bottom block gives the share of the variance of $\bar f$ over all training cell-days carried by each component.}
\label{tab:lgcp_params}
\vspace{2pt}
\small
\setlength{\tabcolsep}{4pt}
\begin{tabular}{llcccc}
\toprule
 & & \multicolumn{2}{c}{Value} & \multicolumn{2}{c}{Effect on $\lambda$} \\
\cmidrule(lr){3-4}\cmidrule(lr){5-6}
Parameter & Description & May & Nov. & May & Nov. \\
\midrule
\multicolumn{6}{l}{\textit{Intensity of active cells} $\log\lambda=\alpha+\mathbf{x}^\top\boldsymbol{\beta}+f$} \\
$\alpha$ & Intercept & $-$0.230 $\pm$ 0.014 & $-$0.212 $\pm$ 0.007 &  &  \\
$\beta_{\mathrm{chl}}$ & Chlorophyll-a (per SD) & 0.114 $\pm$ 0.016 & 0.120 $\pm$ 0.017 & $\times$1.121 & $\times$1.128 \\
$\beta_{\mathrm{temp}}$ & Sea temperature (per $^\circ$C) & 0.107 $\pm$ 0.013 & 0.099 $\pm$ 0.012 & $\times$1.017 & $\times$1.015 \\
$\beta_{\mathrm{sal}}$ & Salinity (per psu) & 0.106 $\pm$ 0.019 & 0.088 $\pm$ 0.021 & $\times$1.088 & $\times$1.073 \\
$\beta_{\mathrm{ban}}$ & Fishing ban (on vs.\ off) & $-$0.660 $\pm$ 0.018 & $-$0.641 $\pm$ 0.013 & $\times$0.101 & $\times$0.151 \\
$\beta_{\mathrm{hol}}$ & Public holiday (on vs.\ off) & $-$0.308 $\pm$ 0.036 & $-$0.273 $\pm$ 0.019 & $\times$0.202 & $\times$0.218 \\
$\beta_{\mathrm{wkend}}$ & Weekend (on vs.\ off) & $-$0.136 $\pm$ 0.027 & $-$0.087 $\pm$ 0.016 & $\times$0.742 & $\times$0.825 \\
$\beta_{\mathrm{lag1}}$ & Count at $t-1$ (per count) & 0.091 $\pm$ 0.005 & 0.093 $\pm$ 0.006 & $\times$1.106 & $\times$1.102 \\
$\beta_{\mathrm{lag7}}$ & Count at $t-7$ (per count) & 0.035 $\pm$ 0.004 & 0.039 $\pm$ 0.003 & $\times$1.040 & $\times$1.041 \\
\midrule
\multicolumn{6}{l}{\textit{Kernel variances and period}} \\
$\sigma^2_{s,\mathrm{broad}}$ & Spatial RBF, $\ell=0.2$ & 1.982 $\pm$ 0.176 & 1.690 $\pm$ 0.152 & & \\
$\sigma^2_{s,\mathrm{loc}}$ & Spatial RBF, $\ell=0.01$ & 0.200 $\pm$ 0.043 & 0.246 $\pm$ 0.044 & & \\
$\sigma^2_{t,\mathrm{RBF}}$ & Temporal RBF, $\ell=0.01$ & 0.208 $\pm$ 0.013 & 0.211 $\pm$ 0.015 & & \\
$\sigma^2_{t,\mathrm{per}}$ & Periodic, $\ell=0.8$ & 0.266 $\pm$ 0.028 & 0.251 $\pm$ 0.026 & & \\
$p$ & Period (days) & 13.96 $\pm$ 0.13 & 13.98 $\pm$ 0.20 & & \\
\midrule
\multicolumn{6}{l}{\textit{Share of the variance of the posterior-mean latent field} $\bar f$} \\
Static & Time-constant cell map & 79.5\% & 84.8\% & & \\
Periodic & Periodic component & 16.3\% & 11.4\% & & \\
Short-term & Temporal RBF component & 0.7\% & 0.5\% & & \\
\bottomrule
\end{tabular}
\end{table}

\paragraph{Covariate effects.}
Every coefficient has the same sign in all 80 runs. Regulatory and calendar factors have the largest effects: during the fishing ban, the intensity of active cells falls to 10--15\% of its value outside the ban, on public holidays to about 20\%, and on weekends to 74--82\%. Recent activity raises the intensity: each additional count observed in the same cell on the previous day multiplies it by 1.10, and each count observed seven days earlier by 1.04. All three environmental covariates have positive effects: chlorophyll-a raises the intensity by 12--13\% per standard deviation, consistent with fishing activity following primary productivity, salinity by 7--9\% per psu, which may partly reflect differences between 2024 and 2025, and sea temperature by 1.5--1.7\% per $^\circ$C. These are conditional associations learned from observational data, not causal effects.

\paragraph{Latent field.}
Splitting the posterior mean $\bar f$ of the latent field into a time-constant, a periodic and a short-term component shows that 80--85\% of its variance is a time-constant map of cell effects; since the spatial lengthscales are shorter than the cell size, the spatial kernel acts mainly as a cell-level effect rather than as a smooth surface. This map ranks the cells by the intensity of their activity (Spearman correlation 0.77--0.87 with the mean positive count of each cell) and is nearly the same in May and November (correlation 0.95). The periodic component, with a learned period of 14 days, accounts for 11--16\% of the variance and is highest on Sundays and Mondays and lowest on Wednesdays and Thursdays, a residual day-of-week pattern of 3--8\% beyond the weekend indicator. The short-term component accounts for less than 1\%, as day-to-day variation is captured mainly by the lagged counts.
